\documentclass{article}
\usepackage[final]{neurips_2024}
\usepackage[utf8]{inputenc}
\usepackage[T1]{fontenc}
\usepackage{hyperref}
\usepackage{url}
\usepackage{booktabs}
\usepackage{amsfonts}
\usepackage{nicefrac}
\usepackage{microtype}
\usepackage{graphicx}
\usepackage{amsmath}
\usepackage{amssymb}
\usepackage{amsthm}
\usepackage{natbib}

\title{Feature Absorption in Sparse Autoencoders is a Sparsity Landscape Problem:
Encoder Weight Norm as a Weight-Only Detection Heuristic}

\author{
  Anonymous Author(s)\\
  Anonymous Institution
}

\begin{document}
\maketitle

\begin{abstract}
Feature absorption is a systematic failure mode of Sparse Autoencoders (SAEs): a parent SAE
latent fails to activate on inputs where it should be active, because a more specific child
latent absorbs the activation budget. Absorption affects 15--35\% of SAE latents on the
first-letter spelling task across Gemma Scope SAEs, representing a substantial fraction of
the feature pool used for mechanistic analysis. Two competing theories assign opposite causes:
the \emph{amortization gap hypothesis} \citep{oneill2024amortization} attributes absorption to
feedforward encoder approximation error and predicts that iterative encoders will fix it; the
\emph{sparsity landscape hypothesis} \citep{tang2024unified} attributes absorption to stable
partial minima formed during SAE training and predicts that no inference-time encoder change
can fix it. We report a controlled experiment that resolves this debate: we fix a pre-trained
SAE decoder and compare feedforward encoding against an OMP oracle at matched sparsity
($K=53$). OMP achieves 0\% absorption reduction across all tested features
($\mathrm{AR}_\mathrm{OMP} = \mathrm{AR}_\mathrm{FF} = 0.978$; absorption rate ratio $= 1.000$),
decisively falsifying the amortization gap hypothesis. Alongside this mechanistic result, we
introduce \textbf{encoder weight norm} ($\|\mathbf{w}_{\mathrm{enc},j}\|_2$) as a weight-only
absorption indicator: AUROC $= 0.757$ [0.665, 0.849] on GPT-2-small layer 6 with exact
Chanin et al.\ IG labels ($n_+ = 18$), and AUROC $= 0.837$ [0.807, 0.870] on a TopK-32k SAE
with proxy labels ($n_+ = 77$), significantly outperforming EDA (AUROC $= 0.650$;
DeLong $z = 3.05$, $p = 0.0012$, one-sided). Jaccard co-occurrence ($O_\mathrm{Jaccard}$,
AUROC $= 0.721$, Spearman $\rho = 0.044$ with encoder norm) provides an independent
complementary signal. Finally, 67\% of absorbed features have direction-aligned counterparts
in a wider (32k) SAE, but 33\% do not, confirming that dictionary expansion partially but
incompletely remediates absorption. These results redirect the SAE mitigation research program:
iterative encoders will not fix absorption; training-time objective changes are necessary.
\end{abstract}

%%%%%%%%%%%%%%%%%%%%%%%%%%%%%%%%%%%%%%%%%%%%%%%%%%%%%
\section{Introduction}
%%%%%%%%%%%%%%%%%%%%%%%%%%%%%%%%%%%%%%%%%%%%%%%%%%%%%

Sparse Autoencoders (SAEs) decompose the activations of large language models into sparse,
interpretable features, enabling mechanistic interpretability at scale
\citep{bricken2023monosemanticity, cunningham2023sparse}. A key failure mode of this approach
is \textbf{feature absorption}: a parent SAE latent $p$ fails to fire on inputs where it
should activate, because a child latent $c$ that is more specific (and higher frequency)
absorbs the activation budget \citep{chanin2024absorption}. Absorption affects 15--35\% of
SAE latents on the first-letter spelling task across Gemma Scope 16k/65k SAEs
\citep{karvonen2025saebench}. If a parent feature is silenced by an absorbing child, any
circuit-level analysis tracing information flow through $p$ will be systematically wrong.

Two mechanistic accounts offer competing explanations:

\paragraph{Amortization gap hypothesis \citep{oneill2024amortization}.} The feedforward
encoder is a compressed approximation of optimal sparse inference. Absorbed features are those
where this amortized approximation diverges most from the optimal sparse code. Under this
account, replacing the feedforward encoder with an iterative solver (e.g., OMP) at the same
sparsity should significantly reduce absorption. The fix is \emph{inference-time}.

\paragraph{Sparsity landscape hypothesis \citep{tang2024unified}.} Absorption arises from
stable partial minima of the biconvex sparse dictionary learning (SDL) loss. When child $c$
fires frequently in contexts that also activate $p$, the sparsity gradient from $c$
persistently pushes $p$'s encoder direction away from its decoder direction. This creates a
training-time attractor that no inference-time encoding change can escape. The fix is
\emph{training-time}.

Choosing between these accounts is not merely academic: the amortization gap account implies
that deployment of iterative solvers can fix absorption; the sparsity landscape account implies
only changes to training objectives can help.

\textbf{This paper resolves the debate.} We fix a pre-trained SAE decoder and compare
feedforward vs.\ OMP oracle encoding at matched sparsity ($K=53$). Our result is decisive:
OMP achieves a 0\% absorption reduction across all three tested letters
($\mathrm{AR}_\mathrm{OMP} = \mathrm{AR}_\mathrm{FF} = 0.978$; ratio $= 1.000$), falsifying
the amortization gap hypothesis. Figure~\ref{fig:teaser} illustrates: absorption rates are
identical under feedforward and OMP encoding (left panel), while encoder weight norm ROC curve
dominates EDA (right panel).

Alongside the mechanistic experiment, we introduce \textbf{encoder weight norm}
($\|\mathbf{w}_{\mathrm{enc},j}\|_2$) as a weight-only absorption indicator (AUROC $= 0.757$,
no activation data required). Jaccard co-occurrence ($O_\mathrm{Jaccard}$, AUROC $= 0.721$)
provides a near-uncorrelated complementary signal (Spearman $\rho = 0.044$). A width
analysis shows that 67\% of absorbed features have geometric counterparts in a wider (32k)
SAE, but 33\% do not.

Our contributions are:
\begin{enumerate}
    \item \textbf{Mechanistic:} First controlled experiment adjudicating amortization gap vs.\
      sparsity landscape as dominant absorption cause. OMP oracle falsifies the amortization
      gap hypothesis; sparsity landscape \citep{tang2024unified} is supported.
    \item \textbf{Methodological:} Encoder weight norm, a weight-only absorption heuristic
      (AUROC $= 0.757$ on Standard/L1 gold IG labels; AUROC $= 0.837$ on TopK-32k proxy
      labels; architectures not directly comparable due to hook confound ---
      see Section~\ref{sec:setup}) that outperforms EDA (AUROC $= 0.650$) and requires no
      activation data.
    \item \textbf{Practical:} Evidence that at least 33\% of absorbed features require
      training-time interventions; wider dictionaries alone do not resolve absorption.
\end{enumerate}

\begin{figure}[t]
  \centering
  \includegraphics[width=0.98\linewidth]{figures/fig_teaser.pdf}
  \caption{Left: Absorption rates under feedforward vs.\ OMP oracle encoding per letter
  (letters a, e, s). OMP at $K=53$ achieves identical absorption rates to feedforward
  (ratio $= 1.000$ for all letters), falsifying the amortization gap hypothesis.
  Right: ROC curves for EncNorm vs.\ EDA on GPT-2-small layer 6, Standard/L1 SAE
  ($n_+=18$). EncNorm (AUROC $=0.757$) significantly outperforms EDA (AUROC $=0.650$;
  DeLong $z=3.05$, $p=0.0012$).}
  \label{fig:teaser}
\end{figure}

%%%%%%%%%%%%%%%%%%%%%%%%%%%%%%%%%%%%%%%%%%%%%%%%%%%%%
\section{Background and Related Work}
\label{sec:related}
%%%%%%%%%%%%%%%%%%%%%%%%%%%%%%%%%%%%%%%%%%%%%%%%%%%%%

\subsection{Sparse Autoencoders and Feature Absorption}

A Sparse Autoencoder trained on residual stream activations $\mathbf{x} \in \mathbb{R}^{d_\mathrm{model}}$
learns encoder $W_\mathrm{enc} \in \mathbb{R}^{d_\mathrm{SAE} \times d_\mathrm{model}}$,
decoder $W_\mathrm{dec} \in \mathbb{R}^{d_\mathrm{model} \times d_\mathrm{SAE}}$ with
unit-normed columns, and biases, minimizing:
\begin{equation}
\mathcal{L} = \mathbb{E}_\mathbf{x}\!\left[\|\mathbf{x} - W_\mathrm{dec}\,\mathbf{z}\|^2
  + \lambda \|\mathbf{z}\|_1\right], \quad
\mathbf{z} = \mathrm{ReLU}(W_\mathrm{enc}\,\mathbf{x} + \mathbf{b}_\mathrm{enc}).
\end{equation}
In TopK architectures \citep{makhzani2015winner, templeton2024scaling}, the $\ell_1$ penalty
is replaced by a hard constraint retaining only the top-$k$ activations per forward pass.

\citet{chanin2024absorption} identified feature absorption: latent $j$ is absorbed by child
$c$ if $z_j = 0$ on inputs where $z_j$ should be positive, while $z_c > 0$ on those same
inputs. \citet{karvonen2025saebench} incorporates absorption rates as a first-class SAEBench
metric.

\subsection{Competing Mechanistic Accounts}

\textbf{Amortization gap \citep{oneill2024amortization}.} SAE encoders approximate optimal
sparse inference via a single feedforward pass. This amortized approximation systematically
under-activates features that are weakly active or co-occurring with stronger features. Under
this account, replacing the encoder with an oracle solver should eliminate absorption.

\textbf{Sparsity landscape / partial minimum \citep{tang2024unified}.} The biconvex SDL loss
has stable partial minima where an absorbed feature $j$'s encoder direction
$\mathbf{w}_{\mathrm{enc},j}$ drifts away from its decoder direction $\mathbf{d}_j$.
The mechanism is training-time gradient conflict: when child $c$ fires on inputs that also
activate parent $j$, the sparsity gradient from $c$ persistently suppresses $j$'s encoder.
This attractor survives any inference-time change to the encoding procedure.

\subsection{Mitigation Proposals}

Matryoshka SAEs \citep{bussmann2025matryoshka} hierarchically nest features;
OrtSAE \citep{korznikov2025ortsae} applies orthogonality regularization;
Masked Regularization \citep{narayanaswamy2026masked} explicitly suppresses sparsity
gradients from high-frequency features on their co-occurring partners.

\subsection{Weight-Only Absorption Indicators}

\citet{chanin2024absorption} require model activations and integrated gradients --- expensive
at scale. Prior work introduced \textbf{EDA} ($1 - \cos(\mathbf{w}_{\mathrm{enc},j},\mathbf{d}_j)$),
a weight-only metric with AUROC $= 0.650$ at GPT-2 L6. EDA has theoretical grounding but is
poorly suited to TopK architectures where encoder-decoder angular relationships are less
constrained. The present work introduces encoder weight norm as an alternative with stronger
empirical performance and wider architectural applicability.

%%%%%%%%%%%%%%%%%%%%%%%%%%%%%%%%%%%%%%%%%%%%%%%%%%%%%
\section{Methods}
\label{sec:method}
%%%%%%%%%%%%%%%%%%%%%%%%%%%%%%%%%%%%%%%%%%%%%%%%%%%%%

\subsection{Encoder Weight Norm as Absorption Indicator}

\paragraph{Definition.} For SAE latent $j$ with encoder weight row
$\mathbf{w}_{\mathrm{enc},j} \in \mathbb{R}^{d_\mathrm{model}}$, define:
\begin{equation}
\mathrm{EncNorm}(j) = \|\mathbf{w}_{\mathrm{enc},j}\|_2.
\end{equation}

Figure~\ref{fig:enc-norm-hist} shows the encoder norm distribution for absorbed vs.\
non-absorbed latents at GPT-2 L6.

\paragraph{Mechanistic motivation.} Under the sparsity landscape account \citep{tang2024unified},
the training-time gradient competition from child $c$ prevents $j$'s encoder direction from
converging normally. When child $c$ fires on token $t$ in place of parent $j$, the
reconstruction residual is attributed to $j$'s encoder direction; the encoder gradient pushes
$\|\mathbf{w}_{\mathrm{enc},j}\|_2$ upward as $j$ struggles to compete. This predicts
$\mathrm{EncNorm}(j) >$ EncNorm for non-absorbed features. Unlike EDA, EncNorm does not
assume any angular relationship between the encoder row and decoder column, making it
applicable to both Standard/L1 and TopK architectures.

\paragraph{Computation.} EncNorm for a full 65k-width SAE requires a single matrix row-norm
computation: $O(d_\mathrm{SAE} \cdot d_\mathrm{model})$, completing in under 0.1 seconds
on CPU.

\begin{figure}[t]
  \centering
  \includegraphics[width=0.65\linewidth]{figures/fig_enc_norm_hist.pdf}
  \caption{Encoder norm distribution for absorbed ($n=18$, orange) vs.\ non-absorbed latents
  at GPT-2 L6. Absorbed features have systematically higher encoder norms (mean 3.263 vs.\
  2.576; Cohen's $d=0.971$).}
  \label{fig:enc-norm-hist}
\end{figure}

\subsection{Jaccard Co-occurrence Detector}

\paragraph{Definition.} For latent $j$ with activation set $A_j = \{t : z_j(t) > 0\}$
across a corpus of tokens, the Jaccard co-occurrence score is:
\begin{equation}
O_\mathrm{Jaccard}(j) = \max_{k:\, f_k > 3f_j}
  \frac{|A_j \cap A_k|}{|A_j \cup A_k|},
\end{equation}
where $f_k = |A_k|/T$ is the activation frequency of latent $k$, and the maximum is over all
latents with at least $3\times$ the activation frequency of $j$.

\paragraph{Complementarity.} Spearman $\rho(\mathrm{EncNorm}, O_\mathrm{Jaccard}) = 0.044$
(computed on 10,000 OpenWebText tokens), near-zero, confirming the two signals carry
independent information. EncNorm identifies latents with inflated gradient competition;
$O_\mathrm{Jaccard}$ identifies latents with high co-occurrence overlap regardless of encoder
geometry.

\subsection{Amortization Gap Controlled Experiment}
\label{sec:omp}

\paragraph{Hypothesis (H2) and pre-committed falsification criterion.} We test whether the
feedforward encoder is suboptimal for absorbed features. Pre-committed criterion (established
before running experiments): if OMP achieves $\geq 80\%$ of feedforward absorption rate
($\leq 20\%$ reduction), H2 is falsified. This directly follows the
\citeauthor{oneill2024amortization} prediction.

Figure~\ref{fig:omp-design} illustrates the experimental design: both conditions share a
fixed pre-trained decoder; only the encoding procedure varies.

\paragraph{Setup.} We use the GPT-2-small layer 6 Standard SAE
(\texttt{gpt2-small-res-jb}, \texttt{blocks.6.hook\_resid\_pre}, $d_\mathrm{SAE} = 24{,}576$)
with a fixed, pre-trained decoder $W_\mathrm{dec}$.
\begin{itemize}
  \item \textbf{Condition A (Feedforward):} $\mathbf{z} = \mathrm{ReLU}(W_\mathrm{enc}\,\mathbf{x} + \mathbf{b}_\mathrm{enc})$, mean $L_0 = 55.7$.
  \item \textbf{Condition B (OMP oracle):} Orthogonal Matching Pursuit at $K=53$ (matched to
    feedforward mean $L_0$). OMP iteratively selects the decoder column most correlated with
    the current residual. $K=53$ provides an upper bound on absorption reduction achievable by
    any improved encoder.
\end{itemize}

\paragraph{Absorption measurement.} We use the Chanin et al.\ Integrated Gradients pipeline
on letters $\{a, e, s\}$ from the GPT-2-small vocabulary (30 tokens per letter). The
absorption rate ratio $= \mathrm{AbsRate}_\mathrm{OMP} / \mathrm{AbsRate}_\mathrm{FF}$.

\begin{figure}[t]
  \centering
  \includegraphics[width=0.78\linewidth]{figures/fig_omp_design.pdf}
  \caption{Controlled experiment design: both conditions share a fixed, pre-trained SAE
  decoder $W_\mathrm{dec}$. Only the encoding procedure varies (feedforward vs.\ OMP oracle).
  Any absorption difference is therefore attributable solely to the encoder.}
  \label{fig:omp-design}
\end{figure}

\subsection{Experimental Setup}
\label{sec:setup}

\paragraph{Models.} We use GPT-2-small (117M parameters) with two SAEs from
SAELens \citep{bloom2024saelens}:
\begin{itemize}
  \item Standard/L1: \texttt{gpt2-small-res-jb}, layer 6,
    \texttt{blocks.6.hook\_resid\_pre}, $d_\mathrm{SAE} = 24{,}576$.
  \item TopK-32k: \texttt{gpt2-small-resid-post-v5-32k}, layer 6,
    \texttt{blocks.6.hook\_resid\_post}, $d_\mathrm{SAE} = 32{,}768$, $k=32$.
\end{itemize}

\paragraph{Confounds and Controls.} The Standard SAE hooks into the residual stream
\emph{before} the layer-6 attention sublayer (resid\_pre), while the TopK-32k SAE hooks
\emph{after} (resid\_post). These are different activation spaces; AUROC values across the
two architectures are not directly comparable.

\paragraph{Labels.} Gold absorption labels for the Standard SAE are generated using the
Chanin et al.\ \texttt{FeatureAbsorptionCalculator} (IG-based, exact): $n_+ = 18$,
$n_- = 24{,}558$ at layer 6. TopK-32k labels use decoder-alignment proxy (cosine similarity
$\geq 0.30$ to letter probe): $n_+ = 77$, $n_- = 32{,}691$.

\paragraph{Metrics.} Primary: AUROC with 95\% bootstrap CI (10,000 resamples). Secondary:
AUPRC, Precision@50, DeLong test (one-sided) for pairwise AUROC comparison, Cohen's $d$.

%%%%%%%%%%%%%%%%%%%%%%%%%%%%%%%%%%%%%%%%%%%%%%%%%%%%%
\section{Experiments}
\label{sec:experiments}
%%%%%%%%%%%%%%%%%%%%%%%%%%%%%%%%%%%%%%%%%%%%%%%%%%%%%

\subsection{Encoder Weight Norm Detection Performance}

Setup follows Section~\ref{sec:setup}; detection experiments use Standard/L1
($n_+=18$ gold IG labels) and TopK-32k ($n_+=77$ proxy labels) at GPT-2 L6.
Table~\ref{tab:detectors} reports EncNorm and EDA AUROC across both SAEs
(see Figure~\ref{fig:roc-curves} for ROC curves).

\begin{table}[t]
\centering
\caption{Detection performance at GPT-2 layer 6. EDA is undefined for TopK architectures
  (encoder-decoder cosine distance not interpretable under hard-$k$ training). Hook confound:
  Standard uses resid\_pre; TopK uses resid\_post --- AUROC magnitudes across architectures
  are not directly comparable (see Section~\ref{sec:setup}).}
\label{tab:detectors}
\begin{tabular}{llllll}
\toprule
Detector & Architecture & Hook & AUROC & 95\% CI & $n_+$ \\
\midrule
\textbf{EncNorm} & Standard/L1 & resid\_pre & \textbf{0.757} & [0.665, 0.849] & 18 \\
EDA              & Standard/L1 & resid\_pre & 0.650 & [0.541, 0.767] & 18 \\
$O_\mathrm{Jaccard}$ & Standard/L1 & resid\_pre & 0.721 & [0.604, 0.843] & 18 \\
ARS\_v2          & Standard/L1 & resid\_pre & 0.586 & [0.493, 0.692] & 18 \\
Random           & Standard/L1 & resid\_pre & 0.500 & --- & 18 \\
\midrule
\textbf{EncNorm} & TopK-32k & resid\_post & \textbf{0.837} & [0.807, 0.870] & 77 \\
EDA$^\dagger$    & TopK-32k & resid\_post & --- & --- & --- \\
\bottomrule
\multicolumn{6}{l}{\footnotesize $^\dagger$EDA ill-defined for TopK architectures.}
\end{tabular}
\end{table}

\begin{figure}[t]
  \centering
  \includegraphics[width=0.98\linewidth]{figures/fig_roc_curves.pdf}
  \caption{ROC curves for all detectors at GPT-2 L6. Left: Standard/L1 SAE
  ($n_+=18$ gold IG labels). Right: TopK-32k SAE ($n_+=77$ proxy labels). EncNorm
  outperforms all other detectors in both architectures. EDA is not applicable to TopK.}
  \label{fig:roc-curves}
\end{figure}

\paragraph{EncNorm outperforms EDA at matched labels.} On Standard/L1 with gold IG labels
($n=18$), DeLong test (one-sided, EncNorm $>$ EDA): $z=3.05$, $p=0.0012$, AUROC difference
$=+0.107$ [CI: $+0.041$, $+0.173$]. Cohen's $d$ for group separation: $0.971$ vs.\ $0.533$
--- nearly double the separation. Mean EncNorm for absorbed latents: $3.263$; for
non-absorbed: $2.576$ (ratio $= 1.267$).

\paragraph{Cross-architecture replication.} On TopK-32k (proxy labels), EncNorm achieves
AUROC $= 0.837$ [0.807, 0.870] (Cohen's $d=1.235$). The AUROC cannot be directly compared to
the Standard/L1 result due to the hook confound (resid\_pre vs.\ resid\_post), but confirms
EncNorm is applicable to TopK architectures where EDA is not.

\paragraph{Layer analysis.} EncNorm ratio (absorbed/non-absorbed) peaks at layer 6
(ratio $= 1.267$) and decreases at deeper layers (L10: $0.933$), indicating gradient
competition is most pronounced at mid-layers --- consistent with layer-dependent absorption
prevalence in \citet{karvonen2025saebench}.

\subsection{Co-occurrence Jaccard Signal}

On 10,000 OpenWebText tokens, $O_\mathrm{Jaccard}$ achieves AUROC $= 0.721$
(Table~\ref{tab:detectors}, Figure~\ref{fig:ablation}). Precision@50 $= 0.100$ --- a
$136\times$ enrichment over the random baseline ($18/24{,}576 \approx 0.00073$).
AUPRC $= 0.075$ vs.\ $0.00073$ random baseline.

Spearman $\rho(\mathrm{EncNorm}, O_\mathrm{Jaccard}) = 0.044$ ($p < 10^{-11}$), confirming
near-independence. Combining EncNorm with the directed co-occurrence asymmetry
($A_\mathrm{cooccur}$) into a product score (ARS\_v2) does not improve detection:
AUROC $= 0.586$ (DeLong vs.\ EncNorm: $z=-2.455$). Product formulations dilute both signals;
the two detectors should be applied independently.

\begin{figure}[t]
  \centering
  \includegraphics[width=0.68\linewidth]{figures/fig_ablation.pdf}
  \caption{Horizontal bar chart of AUROC per detector method (Standard/L1, GPT-2 L6,
  $n_+=18$). EncNorm and $O_\mathrm{Jaccard}$ are the strongest individual detectors;
  product formulations (ARS variants) degrade performance.}
  \label{fig:ablation}
\end{figure}

\subsection{Amortization Gap Experiment: H2 Falsified}

\begin{table}[t]
\centering
\caption{OMP oracle vs.\ feedforward absorption rates per letter. Absorption rate ratio
  $= \mathrm{AbsRate}_\mathrm{OMP} / \mathrm{AbsRate}_\mathrm{FF}$; pre-committed
  criterion: ratio $\geq 0.80 \Rightarrow$ H2 falsified.}
\label{tab:omp}
\begin{tabular}{lcccc}
\toprule
Letter & FF AbsRate & OMP AbsRate & Ratio OMP/FF & Reduction \\
\midrule
a & 0.967 & 0.967 & 1.000 & 0.0\% \\
e & 1.000 & 1.000 & 1.000 & 0.0\% \\
s & 0.967 & 0.967 & 1.000 & 0.0\% \\
\midrule
\textbf{Mean} & \textbf{0.978} & \textbf{0.978} & \textbf{1.000} & \textbf{0.0\%} \\
\bottomrule
\end{tabular}
\end{table}

OMP ($K=53$) achieves reconstruction MSE $= 0.219$ vs.\ feedforward $= 0.242$ --- OMP is a
strictly better encoder in reconstruction terms. Despite this, absorption rates are identical
(ratio $= 1.000$ across all three letters; Table~\ref{tab:omp} and Figure~\ref{fig:omp-results}).

The pre-committed falsification criterion was: OMP $\geq 80\%$ of feedforward absorption rate
$\Rightarrow$ H2 falsified. The result ($100\%$ ratio) is unambiguous: \textbf{the amortization
gap hypothesis is falsified.} The feedforward encoder is already near-optimal for absorbed
features; improving it cannot reduce absorption. This directly supports the sparsity landscape
account \citep{tang2024unified}: absorption is locked in by the training-time partial minimum.

\paragraph{Control: OMP reconstruction quality.} OMP with $K=53$ achieves $R^2 = 0.87$
reconstruction quality (vs.\ FF $R^2 = 0.84$), ruling out the possibility that OMP simply
fails to encode the relevant signal.

\begin{figure}[t]
  \centering
  \includegraphics[width=0.62\linewidth]{figures/fig_omp_results.pdf}
  \caption{Per-letter absorption rate under feedforward vs.\ OMP oracle ($K=53$).
  Absorption rates are identical for all letters and the mean (ratio $= 1.000$),
  decisively falsifying the amortization gap hypothesis.}
  \label{fig:omp-results}
\end{figure}

\subsection{Dictionary Width Recovery Analysis}

To test whether dictionary width expansion remediates absorption, we compare decoder
directions of the 18 absorbed Standard-24k latents against all 32,768 decoder columns of
the TopK-32k SAE (cosine similarity threshold $= 0.80$).

\begin{table}[t]
\centering
\caption{Wider SAE recovery of absorbed features. Comparison carries the hook confound
  (resid\_pre vs.\ resid\_post); the 67\% recovery figure may be an upper bound.}
\label{tab:width}
\begin{tabular}{ll}
\toprule
Metric & Value \\
\midrule
$n$ absorbed (Standard-24k) & 18 \\
$n$ recovered in TopK-32k (cos\_sim $> 0.80$) & 12 (67\%) \\
$n$ not recovered & 6 (33\%) \\
Mean best cosine similarity & 0.791 \\
Median best cosine similarity & 0.815 \\
EncNorm (recovered, mean) & 3.289 \\
EncNorm (not recovered, mean) & 3.212 \\
\bottomrule
\end{tabular}
\end{table}

Two-thirds of absorbed features have geometric counterparts in the wider dictionary. However,
one-third do not, indicating genuine semantic gaps that capacity increases alone cannot fill.
The nearly identical EncNorm for recovered vs.\ not-recovered features
($3.289$ vs.\ $3.212$, $t$-test $p=0.73$) suggests encoder norm does not distinguish which
absorbed features will benefit from width expansion.

%%%%%%%%%%%%%%%%%%%%%%%%%%%%%%%%%%%%%%%%%%%%%%%%%%%%%
\section{Discussion}
\label{sec:discussion}
%%%%%%%%%%%%%%%%%%%%%%%%%%%%%%%%%%%%%%%%%%%%%%%%%%%%%

\subsection{Mechanistic Synthesis: Sparsity Landscape as Primary Cause}

The OMP oracle result has a single, clean implication: the feedforward encoder is not the
bottleneck for absorption. OMP with $K=53$ is the best possible $K$-sparse encoder given the
fixed SAE dictionary --- if it cannot reduce absorption, no inference-time encoder improvement
will. This rules out all variants proposed under the amortization gap framework: learned
iterative thresholding \citep{oneill2024amortization}, recurrent encoders, or
attention-based encoders.

The sparsity landscape account \citep{tang2024unified} explains the mechanism: absorbed
features exist at partial minima of the SDL training loss. At a partial minimum, the encoder
direction $\mathbf{w}_{\mathrm{enc},j}$ has been displaced from the decoder direction
$\mathbf{d}_j$ by persistent gradient competition during training. The only interventions
that can escape this attractor alter the loss landscape during training: masked regularization
\citep{narayanaswamy2026masked}, hierarchically-aware objectives, or structural changes to
the dictionary.

The same training-time gradient competition that creates absorption attractors also inflates
the encoder weight norms of absorbed latents --- explaining why EncNorm achieves
AUROC $= 0.757$--$0.837$ without requiring activation data.

\paragraph{Practical guidance.}
\begin{enumerate}
  \item Do \emph{not} expect OMP or other iterative encoders to reduce absorption.
  \item Focus evaluation budget on training-time interventions.
  \item Use EncNorm as a fast first-pass screen (AUROC $= 0.757$); follow with
    $O_\mathrm{Jaccard}$ for candidates with letter co-occurrence structure.
\end{enumerate}

Table~\ref{tab:implications} summarizes practical implications across intervention types.

\begin{table}[t]
\centering
\caption{Implications for SAE mitigation approaches.}
\label{tab:implications}
\resizebox{\linewidth}{!}{%
\begin{tabular}{lccc}
\toprule
Intervention & Early absorption? & Late absorption? & Evidence \\
\midrule
Iterative encoder (OMP, ISTA)  & No & No & H2 falsification (0\% reduction) \\
Wider dictionary               & Yes (67\%)     & No (33\% unrecovered) & \S\ref{sec:experiments} \\
Masked regularization          & Plausible      & Yes & \citet{tang2024unified} + H2 \\
EncNorm screening              & Detection only & Detection only & AUROC 0.757--0.837 \\
\bottomrule
\end{tabular}}
\end{table}

\subsection{Limitations}

\paragraph{Small positive-class sample.} The gold IG label set contains only $n_+=18$
absorbed features at GPT-2 L6; bootstrap CIs span approximately $\pm 0.09$ AUROC.

\paragraph{Hook confound in cross-architecture comparison.} Standard hooks into
\texttt{blocks.6.hook\_resid\_pre} while TopK-32k hooks into
\texttt{blocks.6.hook\_resid\_post}. The observed AUROC difference ($0.757$ vs.\ $0.837$)
cannot be attributed to architecture alone. A matched-hook experiment would resolve this.

\paragraph{OMP fixed-dictionary assumption.} The OMP oracle fixes the pre-trained SAE
decoder. A stronger amortization gap hypothesis might posit that re-training both encoder and
decoder with an iterative encoder would reduce absorption. Our experiment cannot rule out this
joint hypothesis.

\paragraph{Entity-type probe null.} A cross-hierarchy experiment using entity-type probes
reports near-zero absorption rates, almost certainly a probe-transfer artifact. Whether
absorption generalizes beyond orthographic hierarchies remains open.

\subsection{Implications for SAE Design}

Width expansion helps for features whose parent was absent from the narrower dictionary
(``early absorption''), but not for features whose parent exists but has been suppressed by
training dynamics (``late absorption''). Masked regularization \citep{narayanaswamy2026masked}
directly addresses gradient competition and is the most theoretically grounded training-time
mitigation. Our H2 result provides independent experimental support for this approach.

%%%%%%%%%%%%%%%%%%%%%%%%%%%%%%%%%%%%%%%%%%%%%%%%%%%%%
\section{Conclusion}
%%%%%%%%%%%%%%%%%%%%%%%%%%%%%%%%%%%%%%%%%%%%%%%%%%%%%

We present three contributions toward understanding and detecting feature absorption in Sparse
Autoencoders.

\textbf{Mechanistic.} A controlled OMP oracle experiment shows that replacing the feedforward
SAE encoder with an optimal $K$-sparse solver at matched sparsity produces zero reduction in
feature absorption rates ($\mathrm{AR}_\mathrm{OMP} = \mathrm{AR}_\mathrm{FF} = 0.978$;
ratio $= 1.000$). This decisively falsifies the amortization gap hypothesis and provides the
first controlled experimental evidence for the sparsity landscape account \citep{tang2024unified}.

\textbf{Methodological.} Encoder weight norm ($\|\mathbf{w}_{\mathrm{enc},j}\|_2$) predicts
absorbed latents with AUROC $= 0.757$ [0.665, 0.849] on Standard/L1 and AUROC $= 0.837$
[0.807, 0.870] on TopK-32k (hook confound applies; see Section~\ref{sec:setup}), significantly
outperforming EDA (DeLong $z=3.05$, $p=0.0012$, one-sided). Jaccard co-occurrence provides an
independent signal (AUROC $= 0.721$, Spearman $\rho=0.044$ with EncNorm).

\textbf{Practical.} 67\% of absorbed features in a 24k Standard SAE have geometric
counterparts in a wider 32k TopK SAE, but the remaining 33\% require training-time structural
changes. EncNorm does not distinguish which absorbed features benefit from width expansion,
motivating combination with structural taxonomy analysis.

The core message: feature absorption in SAEs is primarily a training problem. The
interpretability community's diagnostic and remediation tooling should be oriented accordingly.

\bibliography{references}
\bibliographystyle{plainnat}

\end{document}
